\section{Conclusion}
\label{sec:conclusion}

We asked whether an adaptive evaluation protocol, one that treats its own nulls as triggers for
experimental redesign, finds more vulnerabilities and produces better-calibrated conclusions than a
static protocol \emph{at the same compute budget}. We answered it with compute-matched replay on
437k real attack-search iterations, a complete $5\times4\times100$ cost-to-discovery matrix, and a
live grid on three open-weights instruct models under a real alignment intervention.

Adaptive protocols do find substantially more at matched compute, but not for the reason the
hypothesis proposes. The gain comes from searching over \families and from learning across a
portfolio which family works. Using a failing run's internal signal as a redesign trigger is worth
approximately nothing (AUC $\approx 0.5$ at every checkpoint) and costs detections at large budgets.
On the second half of the hypothesis, providing stronger evidence for genuine robustness, adaptive
protocols currently make things worse: their nulls carry nominally 95\% bounds with 63--67\%
coverage, and anytime-valid inference does not repair this because the error is in the estimand, not
the stopping rule.

\para{Recommendations for practice.}
\begin{itemize}[leftmargin=*,itemsep=0pt,topsep=2pt]
    \item Pre-allocate budget across at least three \families rather than deepening one. This single
    change accounts for the entire $0.00 \to 1.00$ effect measured here.
    \item Spend adaptivity at the portfolio level, using Thompson sampling or successive halving
    over methods, where nulls aggregate into real evidence.
    \item Do not instrument within-search stall detection unless you have validated that your
    progress signal is diagnostic on the surviving-null population. Measure its AUC first.
    \item Never publish a binomial confidence interval on a null as a bound on vulnerability. Report
    the design space searched and the ceiling attained; a null is a statement about your design, not
    about the model.
    \item Measure intervention effect sizes across \families too. The same intervention looks inert
    ($0.077 \to 0.037$) or strongly effective ($1.000 \to 0.519$) depending purely on which family
    you probe with.
\end{itemize}

\para{Future work.} Three directions follow directly. First, test whether a \emph{better} within-run
signal, such as true elicitation probability from repeated sampling or a learned probe, is
diagnostic where target-token probability is not; this is the one way the E1 verdict could flip.
Second, build a bound whose estimand is elicitability under a \emph{specified} design space, and
evaluate its coverage against held-out \families. Third, extend the portfolio experiment to methods
with heterogeneous costs, since our unit-cost accounting flatters cheap methods.
